\chapter{Résumé}

Ce projet étudie un système de recommandation basé sur les graphes capable de produire des recommandations d'article à article parmi les articles Wikipédia du domaine de l'informatique. L'objectif principal de l'étude est de construire un système qui, à partir d'un article source donné, puisse classer des articles liés sur les plans sémantique et structurel, puis, si souhaité, présenter ces résultats sous la forme d'un parcours d'apprentissage. Contrairement aux systèmes de recommandation classiques fondés sur le profilage des utilisateurs, cette étude adopte une approche entièrement basée sur les requêtes et fonctionne sur des données ouvertes de Wikipédia.

L'ossature expérimentale du projet repose sur le jeu de données Custom-WikiCS. Ce jeu de données constitue une version enrichie de la structure de graphe Wiki-CS à l'aide des textes réels des articles Wikipédia. Dans une première étape, un graphe de 11 701 nœuds a été préparé ; après suppression des nœuds isolés, un graphe nettoyé composé de 11 367 nœuds actifs et de 291 039 arêtes orientées a été obtenu. Les analyses structurelles ont montré que le réseau possède une distribution de degrés à queue lourde, que les articles connectés présentent une similarité sémantique plus élevée que des paires d'articles choisies aléatoirement, et que des méthodes telles que Louvain et Infomap permettent d'extraire des communautés significatives. Ces résultats confirment que la structure d'hyperliens de Wikipédia constitue un signal utile pour la génération de recommandations.

Trois composantes principales ont été utilisées du côté méthodologique. Premièrement, les textes des articles ont été convertis en représentations vectorielles numériques à l'aide de BAAI/bge-m3. Deuxièmement, des représentations structurelles ont été obtenues à partir du graphe via Node2Vec. Troisièmement, des modèles structurels basés sur les réseaux de neurones sur graphes, tels que GCN et GraphSAGE, ont été entraînés. Au début du projet, les vecteurs de représentations sémantiques et structurelles étaient combinés afin de produire des recommandations ; par la suite, plusieurs modèles ont été comparés de manière systématique. Ainsi, des approches purement sémantiques, purement structurelles, hybrides, fusionnées et fondées sur les réseaux de neurones sur graphes ont été évaluées conjointement.

L'une des contributions importantes du projet réside dans la refonte ultérieure du protocole d'évaluation. Dans la configuration initiale, une approche de séparation des arêtes orientées était utilisée, tandis que les modèles structurels fonctionnaient en pratique sur une structure non orientée. Cette situation créait une incompatibilité méthodologique susceptible d'affecter en particulier l'interprétation des résultats de Node2Vec. Un nouveau protocole d'évaluation adapté à la structure non orientée a donc été mis en place : 216 123 arêtes non orientées ont été réparties entre 172 900 arêtes d'entraînement et 43 223 arêtes de test, tout en évitant la création de nouveaux nœuds isolés, et l'ensemble des principaux modèles structurels a été réévalué dans ce cadre corrigé.

Les résultats mettent en évidence deux perspectives différentes. Dans le protocole d'évaluation initial, le modèle Node2Vec non corrigé semblait très performant ; par exemple, dans les résultats du Report-6, la valeur Hit@10 atteignait jusqu'à 16,06 %. Cependant, dans le protocole d'évaluation corrigé, cette supériorité a disparu et les meilleures performances globales ont été obtenues par les modèles hybrides corrigés. En particulier, le modèle *BGE pooled + Node2Vec (corrected)* a fourni l'un des meilleurs résultats en termes de Hit@10, tandis que le modèle *BGE-M3 + Node2Vec (corrected)* est resté très compétitif sur des métriques secondaires de classement telles que le MRR, le Recall et le MAP. En revanche, pour les métriques de discrimination des liens telles que le ROC-AUC et le PR-AUC, les modèles purement structurels sont restés très performants. Cela montre que le meilleur modèle de recommandation n'est pas nécessairement le meilleur modèle global de prédiction de liens.

L'étude ne s'est pas limitée à l'exactitude du classement ; la qualité des recommandations a également été analysée. Des métriques telles que l'augmentation de popularité, le coefficient de Gini, la diversité agrégée, la nouveauté et la distance intra-liste ont montré que certains modèles, bien qu'offrant des performances proches en termes de précision, pouvaient présenter des comportements de recommandation très différents. Par exemple, le modèle GCN Only, bien qu'assez faible dans les premiers rangs, affichait un profil solide en matière de diversité et de biais de popularité. À l'inverse, BGE-M3 Only présentait une plus forte tendance à recommander des articles populaires. Les modèles hybrides corrigés se sont distingués comme la famille la plus équilibrée entre précision et qualité des recommandations.

Du côté orienté utilisateur, une interface web locale a également été développée. Cette interface propose une recherche par titre, le choix du modèle, l'affichage des 20 premières recommandations, un panneau d'aperçu des articles ainsi qu'une vue optionnelle de parcours d'apprentissage. La couche de parcours d'apprentissage avait d'abord été testée avec un petit modèle local et Ollama, mais ceux-ci se sont révélés insuffisants pour produire de manière stable des sorties structurées en une seule génération. Dans l'implémentation actuelle, bien que le classement et l'accès restent locaux, l'organisation des parcours d'apprentissage repose désormais sur des sorties structurées basées sur Gemini, une validation stricte côté application et des mécanismes heuristiques de secours en cas d'échec. Ainsi, la couche LLM a été traitée comme une couche de présentation ajoutée ultérieurement, distincte de la logique principale de recommandation.

Ce travail propose un cadre expérimental économique et reproductible grâce à l'utilisation de données ouvertes, de mécanismes de mise en cache locale et de matériel grand public. L'absence d'utilisation de données personnelles, l'absence de suivi des utilisateurs et la génération des recommandations à partir d'espaces de similarité pré-calculés offrent des avantages importants en matière d'éthique et de sécurité. En conclusion, ce projet apporte à la fois une contribution méthodologique et pratique à la génération de recommandations assistée par graphes pour les articles Wikipédia, tout en montrant clairement comment la conception du protocole d'évaluation peut modifier l'interprétation des résultats.
